\shorttitle{Attributing Temporary Network Admission in OpenNHP}
\shortauthors{}
\title[mode=title]{Attributing temporary network admission to individual requests in OpenNHP: Method and empirical evaluation}
% Author names, affiliations, correspondence and ORCID are pending author confirmation.
\begin{abstract}
A successful application connection after a knock can rely on an earlier admission rule or a rule written by another request. Distinguishing these explanations matters for incident reconstruction and access reviews, especially when the relevant records are incomplete. We present a two-stage method for attributing temporary admission to an individual request. First, a graph encodes candidate associations between requests and action bundles. Queries that remove or force candidate edges determine whether the target request has a successful action in all or some explanations compatible with the observations. We prove that the queries implement these two semantics under the stated observation model. Second, the method checks the resource state before and after the write, verifies a fresh connection, and rules out competing writes to the same resource. Native OpenNHP runs illustrate three distinct outcomes: the request reaches application processing, a rule write succeeds, and new admission is established. On the 24 sealed instances of the fixed regression suite, global queries resolved seven cases left open by local propagation and agreed with exact enumeration on all 24. The additional verdicts arose where per-request capacity constraints excluded alternative explanations. Under light load, full evidence collection increased knock latency by a median paired difference of 0.620 ms. On a separate set of small sealed graphs, the median of block-mean solver times was below 1 ms.
\end{abstract}
\begin{keywords}
Network-Infrastructure Hiding Protocol \sep temporary network admission \sep execution attribution \sep protocol testing \sep access control
\end{keywords}
